\documentclass[11pt]{letter}
\usepackage[a4paper,margin=2.5cm]{geometry}
\usepackage[T1]{fontenc}
\usepackage{times}

\signature{Daniel Bachrathy\\
Department of Applied Mechanics\\
Faculty of Mechanical Engineering\\
Budapest University of Technology and Economics\\
M\H{u}egyetem rkp.~3., H-1111 Budapest, Hungary\\
bachrathy@mm.bme.hu}

\address{Department of Applied Mechanics\\
Budapest University of Technology and Economics\\
M\H{u}egyetem rkp.~3.\\
H-1111 Budapest, Hungary}

\begin{document}

\begin{letter}{Editors-in-Chief\\
Journal of Vibration and Control\\
SAGE Publications}

\opening{Dear Editors,}

Please consider the enclosed manuscript, \emph{``Solution-Operator
Semi-Discretization: General-Order Multiplication-Free Stability Analysis of
Time-Periodic Delayed Systems and the Accuracy--Sparsity Trade-Off''}, for
publication in the \emph{Journal of Vibration and Control} as an original
research article.

The manuscript is the direct continuation of my recent JVC article on the
multiplication-free semi-discretization method (MFSD), which established the
linear-time-complexity representation of the one-period solution operator of
time-periodic delay differential equations. The present work answers the
natural follow-up question: once the analysis is written in that
multiplication-free form, \emph{which integrator should be placed inside each
step?} I show that the per-step integrator is a free design choice --- any
explicit Runge--Kutta or implicit collocation-type scheme embeds through its
Butcher tableau --- and that the resulting accuracy--sparsity continuum
possesses a quantifiable, problem-dependent CPU-time optimum (a ``sweet
spot'') at moderate order and many steps. A Shannon-type sampling bound turns
this observation into an a~priori recipe for selecting the stage number and
step count, validated on four benchmark systems including machining
applications, and demonstrated end-to-end by high-resolution stability charts
computed in seconds. All methods and every figure are reproducible through
the open-source Julia package \textsf{SOSD.jl}.

The paper is deliberately explicit about its relation to prior art: neither
the multiplication-free assembly (my previous JVC article) nor collocation
itself is claimed as new; the contributions are the general-order embedding,
the fair order-versus-step benchmarking methodology, the measured
accuracy--sparsity trade-off, and the resulting practical guidance. I
believe this continuity makes JVC the natural venue, and the same readership
that received the predecessor article will find the sequel directly useful.

The manuscript is not under consideration elsewhere, and I have no conflicts
of interest to declare. The research was supported by the Hungarian
Scientific Research Fund (NKFIH-152125).

Thank you for your consideration.

\closing{Sincerely,}

\end{letter}
\end{document}
